\documentclass[aspectratio=169,UTF8]{ctexbeamer}
\usetheme{Madrid}
\usecolortheme{default}
\definecolor{DeepBlue}{HTML}{2563EB}
\definecolor{DeepInk}{HTML}{172033}
\definecolor{DeepGreen}{HTML}{059669}
\setbeamercolor{structure}{fg=DeepBlue}
\setbeamercolor{title}{fg=white,bg=DeepBlue}
\setbeamercolor{frametitle}{fg=white,bg=DeepBlue}
\setbeamercolor{alerted text}{fg=DeepGreen}
\setbeamertemplate{navigation symbols}{}
\setbeamertemplate{footline}[frame number]
\usepackage{booktabs}
\usepackage{listings}
\usepackage{xcolor}
\usepackage{hyperref}
\hypersetup{colorlinks=true,urlcolor=DeepBlue}
\lstset{
  basicstyle=\ttfamily\scriptsize,
  keywordstyle=\color{DeepBlue},
  stringstyle=\color{DeepGreen},
  commentstyle=\color{gray},
  showstringspaces=false,
  breaklines=true,
  frame=single,
  columns=fullflexible
}
\newcommand{\api}[1]{\texttt{#1}}
\newcommand{\keypoint}[1]{\alert{#1}}
\AtBeginSection[]{
  \begin{frame}{本节路线}
    \tableofcontents[currentsection]
  \end{frame}
}

\title{第4课：思考模式}
\subtitle{推理强度、reasoning\_content 与多轮规则}
\author{DeepSeek API 全栈课程}
\date{文档快照：2026-11-05}
\begin{document}
\frame{\titlepage}
\begin{frame}{学习目标}
\begin{itemize}
  \item 切换思考/非思考模式并设置推理强度
  \item 分离推理内容与最终回答
  \item 正确处理流式推理和多轮历史
  \item 理解工具存在时的特殊回传要求
\end{itemize}
\end{frame}
\section{控制方式}
\begin{frame}{默认行为}
当前模型默认开启思考模式，默认 effort 为 \api{high}。简单抽取、格式转换、低延迟交互可显式关闭；复杂推理、规划、代码分析可开启。\\[8pt]
思考更强通常意味着更高延迟与更多输出 Token，并非所有任务都更好。
\end{frame}
\begin{frame}{OpenAI Chat 两组参数}
\begin{itemize}
  \item 开关：\api{thinking.type = enabled/disabled}
  \item 强度：\api{reasoning\_effort = low/high/max}
  \item thinking 优先；省略时由 effort 决定
  \item \api{reasoning\_effort=none} 可关闭
\end{itemize}
SDK 中 thinking 通过 \api{extra\_body} 传入。
\end{frame}
\begin{frame}{跨协议映射}
\begin{tabular}{lll}
\toprule
协议 & 开关 & 强度\\
\midrule
Chat & thinking.type & reasoning\_effort\\
Anthropic & thinking.type & output\_config.effort\\
Responses & reasoning.effort & reasoning.effort\\
\bottomrule
\end{tabular}
\vspace{6pt}

minimal$\rightarrow$low；medium/xhigh$\rightarrow$high；ultra$\rightarrow$max。
\end{frame}
\begin{frame}[fragile]{显式开启}
\begin{lstlisting}[language=Python]
response = client.chat.completions.create(
    model="deepseek-flash",
    messages=messages,
    reasoning_effort="high",
    extra_body={"thinking": {"type": "enabled"}},
)
reasoning = response.choices[0].message.reasoning_content
answer = response.choices[0].message.content
\end{lstlisting}
\end{frame}
\section{参数与输出}
\begin{frame}{思考模式下的参数}
\begin{itemize}
  \item temperature、presence/frequency penalty 不生效
  \item Flash：top\_p 有效范围 0.95--1；非思考时恒为 1
  \item Pro：思考时忽略 top\_p；工具启用为 0.95，否则 1
  \item max\_tokens 同时容纳推理与回答，预算不足会截断
\end{itemize}
兼容层可能静默接受但忽略参数，不能仅以“未报错”判断生效。
\end{frame}
\begin{frame}{两条输出通道}
\begin{description}
  \item[reasoning\_content] 模型推理文本
  \item[content] 面向用户的最终回答
\end{description}
产品界面可只展示最终答案；日志与存储应遵循隐私、安全和产品政策。不要把思维链当可验证证据。
\end{frame}
\begin{frame}{流式拼接}
同一流中 delta 可能携带 reasoning 或 content。分别累积，不要用 \api{else} 假定两者永不同时出现；字段为空是正常情况。\\[8pt]
最终块检查 finish reason 与 usage，再将结果标记完成。
\end{frame}
\section{多轮规则}
\begin{frame}{不带 tools 的多轮}
历史 \api{reasoning\_content} 无需回传；即使传入也会被忽略，不拼入上下文。保存 assistant 的最终 content 即可延续普通聊天。\\[8pt]
这减少历史体积，也避免把推理文本误当用户可见事实。
\end{frame}
\begin{frame}{带 tools 的多轮}
只要请求携带 \api{tools}，后续请求必须完整回传每轮 assistant 的 \api{reasoning\_content}，即使该轮没有实际工具调用；否则可能返回 400。\\[8pt]
最安全做法：原样追加 SDK 返回的完整 assistant message。
\end{frame}
\begin{frame}{模式选择策略}
\begin{tabular}{ll}
\toprule
任务 & 建议\\
\midrule
分类、抽取、改写 & disabled / none\\
普通问答 & low 或 high，按评测选择\\
数学、复杂代码、规划 & high / max\\
实时 UI & 流式 + 合理 effort\\
\bottomrule
\end{tabular}
用质量、P95 延迟、Token 与失败率共同评估。
\end{frame}
\begin{frame}{常见错误}
\begin{itemize}
  \item 默认开启却误以为是普通快速模式
  \item 只打印 content，误判流式“长时间无输出”
  \item 为推理设置过小 max\_tokens
  \item 带工具时丢掉 reasoning\_content
  \item 把推理文本视为来源或审计证明
\end{itemize}
\end{frame}
\begin{frame}{本课例程}
\begin{enumerate}
  \item \api{01\_thinking.py}：非流式双通道
  \item \api{02\_thinking\_stream.py}：流式分离推理与回答
  \item \api{03\_thinking\_multi\_turn.py}：不带工具的多轮
\end{enumerate}
第6课再实现带工具的完整回传循环。
\end{frame}
\begin{frame}{官方资料}
\small
\url{https://api-docs.deepseek.com/zh-cn/guides/thinking_mode}\\
\url{https://api-docs.deepseek.com/zh-cn/api/create-chat-completion}
\end{frame}
\end{document}
